\documentclass[12pt]{article}
\usepackage{amsmath, amssymb, amsthm, mathtools}
\usepackage{geometry}
\usepackage{hyperref}
\usepackage{microtype}
\usepackage{tikz-cd}
\geometry{
  top=1.25in,
  bottom=1.25in,
  left=1.25in,
  right=1.25in
}
\title{\textbf{Obstruction Cohomology and the Topology of Failure}\\
\large Constraint Closure Across Differential, Recursive, and Semantic Regimes}
\author{Flyxion}
\date{\today}

\newtheorem{definition}{Definition}
\newtheorem{theorem}{Theorem}
\newtheorem{proposition}{Proposition}
\newtheorem{lemma}{Lemma}
\newtheorem{corollary}{Corollary}
\newtheorem{remark}{Remark}

\begin{document}

\maketitle

\begin{abstract}
We propose a unifying principle---\emph{Katz-Admissibility}---which identifies a single
obstruction class underlying differential propagation, recursive closure, and semantic
consistency. Building on the non-abelian Katz formula established by Barz, we show that
local (Higgs-type) and recursive ($p$-curvature-type) defects are not independent failure
modes but equivalent projections of a common cohomological obstruction, related by a
canonical comparison map $\phi$. The non-abelian connection is reinterpreted as a
\emph{lifting problem} rather than a differential operator: admissibility becomes the
condition that a certain Cartesian square commutes, and failure is the obstruction to
completing that square. This geometric reformulation is then interpreted across three
structurally aligned domains: field-theoretic dynamics in the Relativistic Scalar--Vector
Plenum (RSVP) framework, event-driven computation in Spherepop, and semantic merging in
the TARTAN infrastructure. In each domain the same identity---recursive obstruction equals
the pullback of differential obstruction---governs stability, yielding a general theory of
admissibility as topological agreement. We further introduce the notion of the
\emph{Admissibility Log} as a second active filtration of a system's state-space, prove
that the event log is not merely a record but a constraint structure in its own right, and
establish that singularities across physics, computation, and knowledge systems are
uniformly characterizable as regions of comparison-map divergence.
\end{abstract}

\newpage
\tableofcontents
\newpage

\section{Introduction}

\subsection{From Logical Consistency to Topological Agreement}

Classical systems of verification, whether in mathematics, physics, or computation,
have largely been grounded in the paradigm of logical consistency. A system is deemed
valid if no contradictions arise within a fixed formal language. This paradigm,
however, presupposes that correctness is a property of static configurations rather
than of transformations. In dynamical settings---where states evolve, histories
accumulate, and structures are continuously reconstituted---this assumption becomes
inadequate.

A more appropriate notion of stability in such contexts is not the absence of
contradiction, but the preservation of \emph{admissibility}: the capacity for local
constraints to propagate globally without introducing irreparable discontinuities.
These discontinuities, which we will refer to as \emph{tears}, are not merely logical
failures but topological obstructions to the coherent extension of structure.

The transition from logical consistency to topological agreement thus reframes
correctness as a geometric property. A system is stable not because it satisfies a
collection of axioms, but because its local deformations and global reconstructions
remain compatible under all admissible transformations. This shift aligns naturally
with modern developments in algebraic geometry and higher category theory, where
objects are understood through their deformation spaces and gluing conditions. It also
resonates with a principle that has been implicit in the study of constrained dynamical
systems: that \emph{what fails} is as informative as what holds, and the structure of
failure has its own topology.

A decisive technical advance in this direction is provided by Barz's 2026 paper on
non-abelian $p$-curvature \cite{barz}, which establishes a non-abelian analogue of
Katz's classical formula \cite{katz} relating the Higgs field of the Hodge filtration
to the $p$-curvature of the conjugate filtration. The present paper is organized around
the observation that this identity is not merely a theorem of arithmetic geometry but a
structural law of admissibility that applies uniformly to field theory, computation,
and semantics.

\subsection{The Non-Abelian Connection as a Lifting Problem}

A crucial reorientation that drives the present work concerns what a connection
\emph{is}. Classically, a connection $\nabla : E \to \Omega^1_{S/k} \otimes E$ is a
differential operator---something you apply to sections of a bundle. This operator
picture is useful for computation but obscures the underlying structure.

The sheared de Rham stack $(S/k)_{dR_c}$ provides a geometrization that replaces the
operator picture with a diagram picture. In this setting, a non-abelian connection on a
moduli object $\mathcal{M}$ is not a map but a Cartesian square:
\[
\begin{tikzcd}
\mathcal{M} \arrow[r] \arrow[d] & \mathcal{M}' \arrow[d] \\
S \arrow[r] & (S/k)_{dR_c}
\end{tikzcd}
\]
The connection \emph{is} the commutativity of this square. Admissibility is the
condition that the square can be completed. A tear is the obstruction to its
completion. This reframing is not a matter of language: it changes what questions are
natural to ask. Instead of asking whether a connection satisfies equations, one asks
whether a lifting problem has a solution, and if not, what the obstruction is.

This is the sense in which the present framework constitutes a passage from an
\emph{algebraic} to a \emph{geometric} theory of failure. The obstruction to lifting
is a cohomological class, and the main result is that this class is the same regardless
of whether the lifting is attempted in the differential or the recursive direction.

\subsection{Overview of the Main Result}

The central claim of this paper is that differential propagation and recursive closure
are not independent mechanisms of failure, but rather distinct projections of a single
underlying obstruction class. Formally, we identify two defect operators associated to
a system state,
\[
\delta_{\mathrm{diff}} \quad \text{and} \quad \delta_{\mathrm{rec}},
\]
corresponding respectively to infinitesimal (forward) deformation and recursive
(history-dependent) closure.

The main structural identity, which we call \emph{Katz-Admissibility}, asserts that
these defects are related by a canonical comparison map $\phi$:
\[
[\delta_{\mathrm{rec}}] = \phi^*[\delta_{\mathrm{diff}}].
\]

This identity implies that admissibility can be checked equivalently in either regime.
More strongly, it establishes that all failures of admissibility arise from a single
cohomological obstruction class, whose manifestations differ only by the choice of
filtration. The comparison map $\phi$ is canonical: it is not a choice or a
convention, but the unique map induced by the relationship between the Hodge and
conjugate filtrations on the moduli stack of admissible structures.

\subsection{Domains of Application}

The significance of Katz-Admissibility extends across three distinct but structurally
aligned domains.

In the context of algebraic geometry, it arises from the non-abelian Katz formula,
which relates the Higgs field associated to the Hodge filtration with the $p$-curvature
associated to the conjugate filtration. This provides a rigorous mathematical
foundation for the equivalence of differential and recursive defects.

In the Relativistic Scalar--Vector Plenum (RSVP) framework, this identity admits a
field-theoretic interpretation. Differential curvature and entropic irreversibility
are recast as dual projections of a unified admissibility defect, suggesting a
structural linkage between geometry and thermodynamics that goes beyond analogy.
Gravity and entropy emerge as the differential and recursive appearances, respectively,
of the same underlying obstruction class.

In the Spherepop computational model and the broader Semantic Infrastructure program,
Katz-Admissibility becomes an operational principle. State transitions, event logs,
and semantic merges are validated not by syntactic rules alone, but by the requirement
that their associated obstruction classes agree under the comparison map. This yields a
notion of computation governed by topological consistency rather than purely symbolic
correctness. Crucially, the event log is not a passive record of transitions but an
\emph{active filtration} of the state-space: it constitutes a second coordinate system
on the space of admissible configurations, and Katz-Admissibility is the condition
that these two coordinate systems agree.

Taken together, these perspectives support a unified view in which admissibility is
the fundamental invariant of systems that evolve, remember, and interact.

\section{Geometric Background}

\subsection{Sheared de Rham Stacks}

Let $k$ be a field of characteristic $p > 0$ and $S$ a smooth $k$-scheme.
Classically, a connection on a vector bundle $E$ over $S$ is given by a map
\[
\nabla : E \to \Omega^1_{S/k} \otimes E
\]
satisfying the Leibniz rule. This formulation treats the connection as additional
structure imposed on an otherwise independent geometric object: the bundle comes first,
and the connection is attached to it afterward.

The sheared de Rham stack $(S/k)_{dR_c}$, introduced in the work of Bhatt, Kanaev,
Mathew, Vologodsky, and Zhang \cite{bhattkanaev}, provides a conceptual reorganization
of this notion. Rather than equipping a bundle with a connection, one instead considers
vector bundles \emph{on} $(S/k)_{dR_c}$ directly. A fundamental result is that
\[
\mathrm{Vect}((S/k)_{dR_c}) \simeq \{ \text{vector bundles on $S$ with flat connection} \}.
\]
Thus, the data of a flat connection is encoded geometrically. Infinitesimal motion and
parallel transport are no longer external operations but are built into the structure of
the ambient space. The connection ceases to be an operator and becomes a place.

Intuitively, $(S/k)_{dR_c}$ enlarges $S$ by formally identifying points that are
infinitesimally close in a manner compatible with divided powers. The resulting object
is not a scheme but a stack, reflecting the fact that infinitesimal identifications
carry nontrivial automorphism data. The groupoid structure of the stack records the
higher coherences of infinitesimal motion that a quotient would destroy.

The significance of this for the present framework is the following. The standard
picture of a connection as an operator encourages the question: \emph{does this
operator have the right properties?} The stack picture encourages the question:
\emph{does this lifting problem have a solution?} These are logically equivalent
questions, but they have different obstruction theories. It is the obstruction theory
of the lifting problem that the Katz formula controls.

\subsection{Non-Abelian Moduli of Connections}

Given a smooth proper morphism $f : X \to S$, the classical object of interest is
the relative de Rham cohomology $H^n_{\mathrm{dR}}(X/S)$, which forms a vector
bundle on $S$ equipped with the Gauss--Manin connection. This bundle is abelian: it
is a module over the structure sheaf, and its connection is an ordinary differential
operator.

The non-abelian generalization replaces this linear object with a moduli stack:
\[
\mathcal{M}_{dR}(X/S, n) := \mathrm{Map}_S\big((X/S)_{dR_c}, B\mathrm{SL}_n\big),
\]
whose $T$-points classify rank $n$ vector bundles with flat connection on
$X \times_S T$ together with a trivialization of the determinant. This is no longer a
module but a geometric object whose points are themselves admissible structures.

The shift from vector bundles to stacks reflects a deeper principle. Instead of
tracking a single cohomological invariant, one studies the entire space of admissible
geometric realizations. The system is not described by what it is, but by the
\emph{moduli of all the ways it could be}. The Gauss--Manin connection is
correspondingly reinterpreted as a geometric structure on $\mathcal{M}_{dR}(X/S,n)$,
arising from a universal lifting property: it is the connection that is present because
the moduli problem is defined over a de Rham base.

This is also where the non-abelian character of the theory becomes essential. In the
abelian (linear) case, different choices of connection can be added and compared.
In the non-abelian case, the moduli stack carries a richer structure: distinct
admissible realizations are not comparable by subtraction but only by the existence of
morphisms between them. The obstruction to finding such morphisms is precisely what
the Katz formula measures.

\subsection{$p$-Curvature as Coarse-Grained Obstruction}

In characteristic $p$, flat connections carry an additional invariant known as the
$p$-curvature. For a connection $\nabla$, the $p$-curvature is a morphism
\[
\psi_p : E \to F^*_{\mathrm{abs}}\Omega^1_{S/k} \otimes E,
\]
measuring the deviation of $\nabla^p$ from linearity under Frobenius.

From the perspective of sheared de Rham stacks, the $p$-curvature arises by
restricting the full infinitesimal structure to a coarser equivalence relation, in
which nilpotent divided power data is collapsed to the condition $a^p = 0$. This
procedure may be understood as a \emph{projection}: one passes from the full stack,
which retains all higher groupoid structure, to a quotient that retains only the
Frobenius residue.

The conceptual point is important. The full sheared de Rham stack is a
\emph{high-fidelity} object: it remembers how points are infinitesimally identified at
all orders. The $p$-curvature is the \emph{residual signal} that remains after this
data is compressed to a single recursive invariant. What is retained is precisely the
obstruction to descent along Frobenius---the failure of the connection to close under
$p$-fold iteration. Everything else is lost in the projection.

This is an instance of a general pattern that will recur throughout the paper. A
complex, structured object admits a canonical lossy projection to a simpler invariant.
The projected invariant captures a specific type of failure and discards the rest. The
remarkable content of the Katz formula is that the invariant obtained by projecting in
the differential direction and the invariant obtained by projecting in the recursive
direction are, in fact, the same object seen through different lenses.

The non-abelian generalization replaces the linear morphism $\psi_p$ with a
stack-theoretic construction, but the conceptual role remains the same: it detects the
obstruction to global coherence under repeated application of local transport. What
changes is that the obstruction now lives in a non-abelian cohomology, so it cannot
be described by a single number or a linear map but only by the geometry of the
moduli of liftings.

\section{Dual Filtrations and Obstruction Theory}

\subsection{The Hodge Filtration as Forward Differential Structure}

Let $S$ be a smooth $k$-scheme. The Hodge filtration arises from a deformation of the
notion of flat connection into a one-parameter family of $\lambda$-connections:
\[
\nabla_\lambda : E \to \Omega^1_{S/k} \otimes E,
\qquad
\nabla_\lambda(f s) = f \nabla_\lambda(s) + \lambda\, df \otimes s.
\]
At $\lambda = 1$ one recovers an ordinary connection; at $\lambda = 0$ one obtains a
Higgs field, which is an $\mathcal{O}_S$-linear map. The parameter $\lambda$ therefore
controls how much of the derivation structure of $\nabla$ is retained. Up to rescaling
by units, these are parameterized by $\mathbb{A}^1/\mathbb{G}_m$, and the
\emph{Hodge-filtered sheared de Rham stack} $(S/k)_{dR_c,+} \to \mathbb{A}^1/\mathbb{G}_m$
encodes this family geometrically.

A \emph{Griffiths-filtered} connection is a filtration $F^\bullet E$ satisfying
\[
\nabla(F^i E) \subseteq \Omega^1_{S/k} \otimes F^{i-1}E,
\]
so the connection lowers filtration degree by one. Passing to the associated graded
\[
\mathrm{gr}_F(E) := \bigoplus_i F^i E / F^{i+1}E,
\]
one obtains a first-order invariant measuring the failure of $\nabla$ to preserve the
filtration. This associated graded structure is the Hodge filtration's way of
\emph{exposing its defect}: where the connection fails to stay within a filtration
level, the associated graded records the deviation.

The Hodge filtration is appropriately called the \emph{forward differential structure}
because it captures the infinitesimal, first-order, local behavior of the connection:
how a system changes at the level of its tangent directions. It does not see recursive
history; it sees only the next step.

\subsection{The Conjugate Filtration as Recursive Compression}

In characteristic $p$, there exists a second, intrinsically arithmetic filtration: the
\emph{conjugate filtration}. Its existence and properties depend fundamentally on the
prime $p$ and on Frobenius, which has no analogue in characteristic zero. This
filtration is the machinery by which the system's recursive behavior---its behavior
under $p$-fold iteration---is encoded geometrically.

The conjugate filtration is constructed by viewing $(S/k)_{dR_c}$ as a torsor over the
Frobenius twist $S'$ and then forming a deformation that interpolates between this
torsor and its split form over $\mathbb{A}^1/\mathbb{G}_m$. Concretely, the
\emph{conjugate-filtered de Rham stack} $(S/k)_{dR_c,c} \to \mathbb{A}^1/\mathbb{G}_m$
encodes a family of structures whose associated graded, at $\lambda = 0$, is a split
torsor over $S' \times B\mathbb{G}_m$ with tangent bundle twist $TS'/k(+1)$. The
filtration is increasing, and it reflects the behavior of the system under Frobenius
iteration.

The key structural observation is that Frobenius is not merely a technical device of
characteristic-$p$ geometry. It is an operation of \emph{recursive compression}: it
takes the full infinitesimal neighborhood of a point and projects it through a
$p$-fold iteration that collapses all intermediate structure. The conjugate filtration
is the filtration that results from allowing this recursive compression to interact with
the geometry of the de Rham stack. Where the Hodge filtration sees the first step, the
conjugate filtration sees the $p$-th.

\subsection{Higgs Field and $p$-Curvature as Obstruction Invariants}

The passage to associated gradeds extracts obstruction data from each filtration.

From the Hodge filtration, one obtains the \emph{Higgs field}
\[
\Theta : \mathrm{gr}_F(E) \to \Omega^1_{S/k} \otimes \mathrm{gr}_F(E)(-1),
\]
which is $\mathcal{O}_S$-linear and measures the deviation of the connection from
preserving the filtration. In the non-abelian setting, this generalizes to a morphism
of stacks
\[
\Theta_{X/S} : \pi^* T_{S/k}(-1) \to T_{\mathcal{M}_{Dol}/(S \times B\mathbb{G}_m)}.
\]
This is the \emph{differential defect}: the amount by which first-order transport fails
to be filtration-preserving.

From the conjugate filtration, one obtains the \emph{conjugate $p$-curvature}
\[
\psi : E \to F^*_{S/k}\Omega^1_{S'/k} \otimes E(+1),
\]
which measures the failure of recursive closure under Frobenius. In the non-abelian
setting, this becomes a morphism of stacks:
\[
\psi_{X/S} : \pi^* F^*_{S/k}T_{S'/k}(+1) \to T_{\mathcal{M}_{Dol,c}/(S \times B\mathbb{G}_m)}.
\]
This is the \emph{recursive defect}: the amount by which $p$-fold iteration of
transport fails to be trivial.

A priori, these are entirely different invariants arising from entirely different
constructions. The Higgs field lives on the Dolbeault moduli stack; the $p$-curvature
lives on the conjugate-filtered Dolbeault moduli stack. They are extracted by different
projection procedures from different filtrations. There is no obvious reason why they
should be related. Yet the Katz formula asserts that they are not merely related but
\emph{identical} after the appropriate transport.

\section{The Non-Abelian Katz Formula}

\subsection{Statement of the Theorem}

Let $f : X \to S$ be a smooth proper morphism of smooth $k$-schemes, and let
$\mathcal{M}_{dR}(X/S,n)$ denote the non-abelian de Rham moduli stack of rank $n$
bundles with flat connection. From the Hodge filtration, one obtains a non-abelian
Higgs field
\[
\Theta_{X/S} : \pi^* T_{S/k}(-1) \to T_{\mathcal{M}_{Dol}/(S \times B\mathbb{G}_m)}.
\]
From the conjugate filtration, one obtains a non-abelian $p$-curvature
\[
\psi_{X/S} : \pi^* F^*_{S/k}T_{S'/k}(+1) \to T_{\mathcal{M}_{Dol,c}/(S \times B\mathbb{G}_m)}.
\]

\begin{theorem}[Non-Abelian Katz Formula, Barz \cite{barz}]
There exists a canonical morphism
\[
\phi : \mathcal{M}_{Dol,c} \to \mathcal{M}_{Dol}
\]
such that
\[
\psi_{X/S} = \phi^*\Theta_{X/S}.
\]
More precisely, there is a Cartesian diagram relating $\mathcal{M}_{Dol,c}$ and
$\mathcal{M}_{Dol}$, and under this diagram the pullback of the non-abelian Higgs
field $\Theta_{X/S}$ equals the conjugate-filtered $p$-curvature $\psi_{X/S}$.
\end{theorem}

This identity is to be understood as an equality of morphisms after pullback along
$\phi$, and it is an equality of stack-valued morphisms, not merely of linear maps.
The Cartesian diagram structure is essential: $\phi$ is not an arbitrary morphism but
one whose definition is forced by the relationship between the two filtrations.

\subsection{Conceptual Interpretation: Two Projections of One Defect}

The non-abelian Katz formula establishes that the two filtrations introduced in the
previous section are not independent structures. Their associated obstructions are
equivalent under the canonical comparison $\phi$.

To understand why this is nontrivial, observe the following. The Higgs field
$\Theta_{X/S}$ arises from the failure of the connection to preserve the Hodge
filtration. It is a \emph{forward} invariant: it measures what happens when you try to
extend the connection one step in the differential direction and find that the
filtration degree drops. The $p$-curvature $\psi_{X/S}$ arises from the failure of
recursive closure under Frobenius. It is a \emph{recursive} invariant: it measures
what happens when you apply the connection $p$ times and find that the result is not
what flatness would predict.

These are measurements taken at different scales and in different directions. The Higgs
field measures a deviation at the tangent level; the $p$-curvature measures a deviation
at the level of a global, $p$-fold operation. The identity $\psi_{X/S} = \phi^*\Theta_{X/S}$
asserts that despite this apparent difference, the obstructions they detect are
identical. The recursive failure is the pullback of the differential failure.

Put differently: if you know where the system tears differentially, you know exactly
where it will tear recursively, and vice versa. The two failure modes are not
independent sources of information about the system. They are two coordinate
representations of a single underlying fact.

\subsection{The No Independent Tear Principle}

The conceptual content of the Katz formula can be distilled into a principle that
applies well beyond its algebraic-geometric origin.

\begin{theorem}[No Independent Tear Principle]
Let $\mathcal{F}$ be a geometric or computational system equipped with both a
differential and a recursive structure. Then all admissibility failures in
$\mathcal{F}$ arise from a single cohomological obstruction class. Equivalently, the
defects detected by differential propagation and recursive closure are equivalent under
a canonical comparison map.
\end{theorem}

\begin{proof}[Sketch]
The non-abelian Katz formula provides an explicit identification of the Higgs field
and $p$-curvature via pullback along $\phi$. Both $\Theta$ and $\psi$ arise as
associated graded invariants of filtrations on the same underlying structure, namely
the moduli stack $\mathcal{M}_{dR}(X/S,n)$ equipped with its two natural filtrations.
Since these filtrations are defined on a common geometric object and their associated
graded pieces are related by the Cartesian diagram of the Katz formula, the induced
obstruction classes must coincide in cohomology. Thus any failure detected in one
regime is necessarily present in the other, and no failure exists in one that is absent
in the other.
\end{proof}

\begin{corollary}
Admissibility is checkable from either regime without loss of information. A system
that is admissible differentially is admissible recursively, and conversely.
\end{corollary}

The No Independent Tear Principle has a further consequence that is worth making
explicit. Since all admissibility failures are projections of a single class, the
structure of \emph{failure space} is simpler than the structure of \emph{state space}.
States can be arbitrarily complex; failures all live in $H^1(\mathcal{F})$. This is a
significant compression, and it is the compression that makes a unified theory of
admissibility possible.

\subsection{The Comparison Map $\phi$ as a Transport Equivalence}

The comparison map $\phi : \mathcal{M}_{Dol,c} \to \mathcal{M}_{Dol}$ has an
interpretation that deserves emphasis. It is not an arbitrary map between two
unrelated moduli spaces. It is the map that aligns the two filtrations---it is, in a
precise sense, the act of choosing a coordinate system on the space of admissible
structures.

In more abstract terms, $\phi$ plays the role of what might be called a
\emph{transport equivalence}: a canonical identification between two representations
of the same structure that are a priori distinct. The Hodge filtration gives one
representation; the conjugate filtration gives another. The map $\phi$ is the proof
that these representations are equivalent, and the Katz formula is the statement that
under this equivalence, the two obstruction invariants coincide.

This can be expressed in the RSVP framework as follows. Define two operators on the
space of field configurations:
\[
\mathcal{J}_{\mathrm{diff}} : X \mapsto \delta_{\mathrm{diff}}(X),
\qquad
\mathcal{R}_{\mathrm{frob}} : X \mapsto \delta_{\mathrm{rec}}(X).
\]
The transport equivalence asserts that there exists a canonical isomorphism
$\mathcal{T}$ of the state-space such that
\[
\mathcal{R}_{\mathrm{frob}}(X) = \mathcal{T}^{-1} \circ \mathcal{J}_{\mathrm{diff}} \circ \mathcal{T}(X).
\]
Recursive dynamics is therefore not new dynamics. It is differential dynamics seen
through a different coordinate system on the admissibility manifold. The choice between
analyzing a system differentially and analyzing it recursively is a choice of
coordinates, not a choice of subject matter.

\section{RSVP Field Interpretation}

\subsection{Field-Space Completion}

In the Relativistic Scalar--Vector Plenum (RSVP) framework, a physical or
computational state is modeled as a triple
\[
X = (\Phi, \vec{v}, S),
\]
where $\Phi$ is a scalar potential, $\vec{v}$ is a vector flow field, and $S$ is an
entropy-like functional encoding structural irreversibility.

The passage from a base space $S$ to its sheared de Rham stack $(S/k)_{dR_c}$ admits a
natural interpretation in this setting as a \emph{field-space completion}. Rather than
treating evolution as an operator acting on $X$, one considers the prestack of all
infinitesimal continuations of $X$ consistent with its internal constraints. This is
the space of all admissible next states, together with the structured ways they relate
to one another. A connection in this context is a rule for lifting trajectories into
this completed space---a specification of which continuations are admissible and how
they fit together.

The field-space completion replaces the question ``what does this state evolve into?''
with the question ``what is the space of admissible evolutions of this state?'' The
difference is not merely philosophical. In the former picture, the dynamics is a
function; in the latter, it is a section of a bundle over the admissibility manifold.
Admissibility failures become visible as the inability to find such a section.

\subsection{Differential vs Recursive Dynamics}

Two distinct but structurally related modes of evolution arise within RSVP. The
differential mode, governed by
\[
\mathcal{J}_{\mathrm{diff}} : X \mapsto \delta_{\mathrm{diff}}(X),
\]
captures infinitesimal change in the field configuration. It corresponds to the Hodge
filtration in the geometric setting and operates at the level of first-order
perturbations of the field. The recursive mode, governed by the Frobenius-type
operator
\[
\mathcal{R}_{\mathrm{frob}} : X \mapsto \delta_{\mathrm{rec}}(X),
\]
captures the effect of repeated or accumulated transformations. It corresponds to the
conjugate filtration and operates at the level of the field's long-run behavior under
iterated application.

These two operators are not independent. By the transport equivalence established in
the previous section, they are related by
\[
\mathcal{R}_{\mathrm{frob}}(X) = \mathcal{T}^{-1} \circ \mathcal{J}_{\mathrm{diff}} \circ \mathcal{T}(X),
\]
where $\mathcal{T}$ is the canonical coordinate change between the Hodge and conjugate
representations. The local differential behavior and the global recursive behavior are
therefore not two aspects of the field that must be separately managed; they are one
aspect expressed in two equivalent notations.

\subsection{Gravity as Differential Obstruction}

Within RSVP, curvature is understood as the failure of local field configurations to
extend consistently across space. This is precisely the role of the Higgs field
$\Theta_{X/S}$: it measures the deviation of a connection from preserving the Hodge
filtration, which is to say, the degree to which local differential data fails to be
globally coherent.

We therefore interpret
\[
\mathcal{G}(X) \sim \Theta_{X/S}
\]
as the \emph{gravitational field}: a differential obstruction encoding the local
curvature of the plenum. In this interpretation, the geometry of space is not a
background against which dynamics occurs but a record of where differential
admissibility fails. Curvature is the map of failure at the infinitesimal scale.

Mass and energy, in this picture, are the sources of the Hodge defect: they are
the field configurations that make it impossible for the differential filtration to be
preserved. The gradient of the gravitational potential is therefore not a force but an
admissibility gradient---a measure of how rapidly the obstruction class changes as one
moves through field-space.

\subsection{Entropy as Recursive Obstruction}

Entropy is associated with the accumulation of irreversible transformations over time.
The $p$-curvature $\psi_{X/S}$ captures precisely such recursive deviation, measuring
the failure of a system to close under repeated application of its dynamics. We
therefore interpret
\[
\mathcal{E}(X) \sim \psi_{X/S}
\]
as the \emph{entropic field}: a recursive obstruction encoding the historical structure
of the system.

Entropy in this framework is not a measure of microscopic disorder in the statistical
sense. It is a geometric trace of the accumulated failure of recursive closure. Every
irreversible transformation contributes to the $p$-curvature of the system's trajectory
in field-space, and the total entropy is the integrated magnitude of this curvature
over the system's history.

The arrow of time, correspondingly, is not an external condition imposed on the
dynamics but a consequence of the structure of the recursive obstruction. Time flows in
the direction of increasing $\mathcal{E}(X)$, which is to say, in the direction in
which recursive closure becomes progressively harder to achieve. The system cannot
``go back'' not because of a microscopic irreversibility but because the obstruction
class accumulated in $\psi_{X/S}$ cannot be unwound by local differential operations.

\subsection{The Linkage Equation and Its Consequences}

The non-abelian Katz formula implies the following identity in the RSVP setting:
\[
\mathcal{E}(X) = \phi^* \mathcal{G}(X),
\]
where $\phi$ is the canonical comparison map between recursive and differential regimes.
This \emph{linkage equation} expresses a fundamental equivalence: the entropic
accumulation of a system is the pullback of its local curvature under recursive
closure.

The consequences of the linkage equation are substantial. First, gravity and entropy
are not independent phenomena; they are dual manifestations of a single admissibility
defect. A system that is locally curved is necessarily historically irreversible, and
one that has accumulated entropy is necessarily inhabiting a curved region of
field-space. Second, singularities admit a unified characterization: a singularity is a
region of field-space where the comparison map $\phi$ fails to admit a consistent
pullback. In classical general relativity, a singularity is a point of divergent
curvature; in the present framework, it is a region where the admissibility defect
diverges so sharply that no admissible state can exist there. This reframing suggests
that singularity resolution is not a matter of regularizing divergent quantities but of
constructing an extension of $\phi$ across the region of failure.

Third, the linkage equation implies that the field equations of RSVP can be stated in
a single obstruction-theoretic form rather than as two separate field equations for
$\mathcal{G}$ and $\mathcal{E}$. This is the sense in which the present framework
``unifies'' gravity and entropy: not by deriving one from the other but by identifying
them as projections of one thing.

\section{Spherepop: Katz-Admissibility as a Type Rule}

\subsection{States, Events, and Transitions as Lifting Problems}

Let $\sigma$ denote a system state in a Spherepop program, and let
\[
\Delta : \sigma \to \sigma'
\]
be an event proposing a transition. In classical operational semantics, such a
transition is validated by checking whether $\Delta$ has the right type and whether
$\sigma$ satisfies the preconditions of $\Delta$. This is a purely local check: it
examines the immediate configuration without reference to history.

In the Katz-admissibility framework, we interpret $\Delta$ differently: as a
\emph{lifting problem}. A transition proposes that the system, currently in state
$\sigma$, can extend to a new state $\sigma'$ through an admissible deformation. This
is a request to complete a diagram:
\[
\begin{tikzcd}
\sigma \arrow[r, dashed, "\Delta"] \arrow[d] & \sigma' \arrow[d] \\
\mathcal{A} \arrow[r, "\phi"] & \mathcal{A}'
\end{tikzcd}
\]
where $\mathcal{A}$ and $\mathcal{A}'$ are the admissibility structures of the current
and proposed states. The transition is admissible if and only if this diagram can be
completed in a way that is compatible with both the differential structure of the state
and the recursive structure of the event log.

This reframing is not cosmetic. By treating transitions as lifting problems, we gain
access to the full machinery of obstruction theory. The question is no longer
``does $\Delta$ type-check?'' but ``does $\Delta$'s lifting problem have a solution,
and if not, what is the obstruction?''

\subsection{The Admissibility Log as Second Filtration}

A central claim of this framework is that the event log of a Spherepop program is not
a passive record but an \emph{active filtration} of the state-space. This deserves
careful statement.

\begin{definition}[Admissibility Log]
Let $\sigma_0, \sigma_1, \ldots, \sigma_n$ be a sequence of states produced by events
$\Delta_1, \ldots, \Delta_n$. The \emph{Admissibility Log} $\mathrm{AL}(\sigma_\bullet)$
is the decreasing filtration
\[
\mathrm{AL}^k(\sigma_\bullet) := \{ \text{admissible states reachable from } \sigma_{n-k}
\text{ without violating the history of } \Delta_1, \ldots, \Delta_{n-k} \}
\]
for $k = 0, 1, \ldots, n$.
\end{definition}

The Admissibility Log is a filtration of the reachable state-space by the depth of
historical constraint. At level $k = 0$, only the current state and its immediate
continuations are visible. At level $k = n$, the entire history imposes its
constraints. The associated graded pieces of this filtration measure how much of the
state-space is ruled out by each additional step of history.

The crucial observation is that this log-filtration is exactly the conjugate filtration
in the geometric sense. The event log is the Frobenius structure of the computation:
it is the mechanism by which the system's past is recursively compressed into a
constraint on its future. A computation that ignores its event log is like a connection
with trivial $p$-curvature---it may be locally flat, but it is making a claim about
global closure that it has not verified.

\subsection{Differential and Recursive Defects}

To each proposed transition $\Delta : \sigma \to \sigma'$, we associate two defect
operators.

The \emph{differential defect} $\delta_{\mathrm{diff}}(\sigma,\Delta)$ measures the
failure of the local state update to preserve admissibility under infinitesimal
deformation. It is computed from the immediate effect of $\Delta$ on the structure of
$\sigma$: does the transition preserve the system's local invariants, its type
constraints, its resource bounds? The differential defect is what classical type-checking
computes.

The \emph{recursive defect} $\delta_{\mathrm{rec}}(\sigma,\Delta)$ measures the
failure of the system to remain admissible when the transition is incorporated into the
Admissibility Log and evaluated under recursive closure. It captures inconsistencies
that emerge only through accumulation: does adding $\Delta$ to the event log create a
historical configuration that cannot be consistently extended? Does it produce a $k$-th
level admissibility violation that the $(k-1)$-th level check missed?

The recursive defect is what distinguishes Katz-admissibility from ordinary
type-checking. It is the obstruction that classical semantics cannot see, because
classical semantics does not treat the event log as a filtration.

\subsection{The Katz-Admissibility Type Rule}

With these definitions in place, the Katz-admissibility condition for transitions is
precisely stated.

\begin{definition}[Katz-Admissibility of Transitions]
A transition $\Delta : \sigma \to \sigma'$ is \emph{Katz-admissible} if and only if
\[
[\delta_{\mathrm{rec}}(\sigma,\Delta)] = \phi^*[\delta_{\mathrm{diff}}(\sigma,\Delta)],
\]
where $\phi$ is the canonical comparison map from the recursive graded to the
differential graded of the admissibility filtrations, and the brackets denote
cohomology classes in $H^1(\mathcal{F}_\sigma)$.
\end{definition}

This yields the following pair of type-inference rules. The admissibility rule is:
\[
\frac{
\mathrm{Check}_{Katz}(\sigma,\Delta) = \mathrm{true}
}{
\sigma \xrightarrow{\Delta} \sigma' : \mathrm{Admissible}
}
\]
and the tear rule is:
\[
\frac{
\mathrm{Check}_{Katz}(\sigma,\Delta) = \mathrm{false}
}{
\sigma \xrightarrow{\Delta} \bot_{\mathrm{NAT}}
}.
\]
The comparison function is:
\[
\mathrm{Check}_{Katz}(\sigma,\Delta)
=
\Big([\delta_{\mathrm{rec}}(\sigma,\Delta)] = \phi^*[\delta_{\mathrm{diff}}(\sigma,\Delta)]\Big).
\]

\begin{definition}[Non-Abelian Tear]
The failure state $\bot_{\mathrm{NAT}}$ is called a \emph{Non-Abelian Tear}. It
represents the topological obstruction to completing the lifting diagram associated to
$\Delta$. A Non-Abelian Tear is not a conventional runtime error; it is a failure of
the system's local and global admissibility structures to agree.
\end{definition}

\subsection{The Non-Abelian Tear as Singularity}

The Non-Abelian Tear deserves to be understood as the computational analogue of a
physical singularity. In the RSVP setting, a singularity was characterized as a region
where the comparison map $\phi$ fails to admit a consistent pullback. In the
computational setting, a Non-Abelian Tear is a transition where the comparison map
between the recursive and differential filtrations of the Admissibility Log cannot be
completed.

This analogy is more than metaphor. In both cases, the failure is the inability to
extend a local admissibility structure to a global one. In both cases, the obstruction
lives in $H^1(\mathcal{F})$ for an appropriate sheaf $\mathcal{F}$. And in both
cases, the appropriate response is not to patch the immediate failure but to understand
the global topology of the admissibility manifold that produced it.

A Non-Abelian Tear signals that the computation has reached a state from which no
admissible continuation exists given the system's history. The state itself may be
locally valid---it may pass all local type-checks---but its position in the
Admissibility Log is incoherent with its differential structure. The computation must
either be rolled back to a state from which an admissible path exists, or the system
must be redesigned so that the offending region of the admissibility manifold is
avoided.

\subsection{Implementation: Practical Computation of $\phi$}

In a practical implementation, the comparison map $\phi$ must be computed from the
data of the transition. A pragmatic definition proceeds as follows.

Let $\mathrm{LocalGr}(\sigma, \Delta)$ denote the associated graded of the
differential filtration induced by the proposed local mutation of $\sigma$ under
$\Delta$. Let $\mathrm{HistoryGr}(\sigma, \Delta)$ denote the associated graded of the
recursive filtration induced by appending $\Delta$ to the Admissibility Log.

The comparison map is then:
\[
\phi : \mathrm{HistoryGr}(\sigma,\Delta) \to \mathrm{LocalGr}(\sigma,\Delta),
\]
and $\mathrm{Check}_{Katz}(\sigma,\Delta)$ tests whether the obstruction classes
induced by these two associated gradeds coincide under $\phi$.

In the simplest implementations, this check can be approximated by a hash of the
differential type signature of $\Delta$ against a hash of the relevant suffix of the
Admissibility Log. Exact computation requires maintaining a richer representation of
the filtration structure, but even the approximation captures the essential
distinction: a transition is checked not just for local type correctness but for
global historical coherence.

\section{Semantic Infrastructure and Sheaf Merge}

\subsection{Knowledge as Sheaves over Contexts}

We model semantic objects---documents, codebases, theories, knowledge graphs---as
sections of a sheaf $\mathcal{K}$ over a base of contexts $U \subseteq \Omega$. A
section $s \in \mathcal{K}(U)$ represents a locally consistent body of knowledge over
the domain $U$, together with restriction maps satisfying the standard gluing axioms
of sheaf theory.

In this setting, consistency is not absolute but \emph{contextual}: two sections may
agree on their overlapping domains while differing on regions outside the overlap.
The problem of integrating multiple knowledge sources is therefore a problem of
sheaf-theoretic descent. One asks: given sections $A \in \mathcal{K}(U)$ and
$B \in \mathcal{K}(V)$ that agree on the overlap $U \cap V$, can one find a section
$C \in \mathcal{K}(U \cup V)$ that restricts to both?

The obstruction to such a global section is measured by the first \v{C}ech cohomology
$H^1(U \cap V, \mathcal{K})$. When this group is trivial, descent succeeds and the
merge exists. When it is nontrivial, the sections cannot be glued without
introducing a tear.

\subsection{The Two Filtrations of a Semantic Field}

In practice, semantic systems carry structure beyond their content. Each section has a
\emph{derivational history}: a record of the inferences, transformations, and
compositions that produced it. This history is not merely metadata; it is part of the
semantic content of the section, because it determines which further inferences are
valid and which connections to other sections are coherent.

This observation motivates two natural filtrations on the sheaf $\mathcal{K}$.

The \emph{differential filtration} $F^\bullet_{\mathrm{diff}} \mathcal{K}$ organizes
sections by the depth of local inference: sections at level $i$ are those whose
content can be reached from the base context by $i$ local rewrite or inference steps.
The associated graded $\mathrm{gr}^i_{\mathrm{diff}} \mathcal{K}$ captures the content
that is introduced at exactly inference depth $i$, which is to say, the content that
requires exactly $i$ steps of local reasoning to reach.

The \emph{recursive filtration} $F^\bullet_{\mathrm{rec}} \mathcal{K}$ organizes
sections by the depth of provenance: sections at level $j$ are those whose derivational
history can be traced to exactly $j$ ancestral sections. The associated graded
$\mathrm{gr}^j_{\mathrm{rec}} \mathcal{K}$ captures the content that is introduced by
exactly the $j$-th level of compositional provenance.

These two filtrations are in general distinct: a section can be locally simple (low
differential level) while having a complex derivational history (high recursive level),
or vice versa. Katz-Admissibility is the condition that governs when they are
compatible.

\subsection{Katz-Admissibility for Merge Operations}

Given two semantic fields $A \in \mathcal{K}(U)$ and $B \in \mathcal{K}(V)$, define
the \emph{differential incompatibility} $\delta_{\mathrm{diff}}(A,B)$ as the
obstruction class arising from comparing their local inference structures on
$U \cap V$---the degree to which their rewrite rules, constraint systems, and
immediate content disagree on the overlap. Define the \emph{recursive incompatibility}
$\delta_{\mathrm{rec}}(A,B)$ as the obstruction class arising from comparing their
provenance graphs on $U \cap V$---the degree to which their derivational histories,
compositional origins, and ancestral dependencies disagree.

\begin{definition}[Katz-Admissibility of Merges]
A merge $A \sqcup B$ exists as a globally coherent section in $\mathcal{K}(U \cup V)$
if and only if
\[
[\delta_{\mathrm{rec}}(A,B)] = \phi^*[\delta_{\mathrm{diff}}(A,B)],
\]
where $\phi$ is the canonical comparison map from the recursive graded to the
differential graded of $\mathcal{K}$ restricted to $U \cap V$.
\end{definition}

This condition is the semantic analogue of the non-abelian Katz formula. The merge
succeeds not merely when the two sections agree on their immediate content (which
would be the condition for local consistency) but when their local content-disagreement
and their historical-provenance-disagreement are related by the canonical comparison.

The significance of this condition is that it captures a type of incompatibility that
purely local checks miss. Two knowledge graphs can have perfectly compatible local
content---the same entities, the same relation types, the same immediate inferences---
while being historically incompatible: they were derived by different paths from
conflicting assumptions, and merging them would create a system that locally appears
consistent but globally is not.

\subsection{Consistency Without Convergence}

A crucial consequence of Katz-admissibility for merges is that global convergence to a
single canonical form is not required. Two semantic fields may remain distinct yet
admissible, provided their obstruction classes coincide under the comparison map.

This yields a precise formulation of what was previously called \emph{consistency
without convergence}. Two sections $A$ and $B$ are \emph{Katz-consistent} if
$[\delta_{\mathrm{rec}}(A,B)] = \phi^*[\delta_{\mathrm{diff}}(A,B)]$, regardless of
whether there exists a third section $C$ to which both restrict. Katz-consistency is a
relation between sections, not a property of individual sections, and it does not
require either section to change.

\begin{proposition}
Katz-consistency is symmetric and, under mild regularity conditions on the sheaf
$\mathcal{K}$, transitive. It defines an equivalence relation on the set of
locally consistent sections, coarser than isomorphism but finer than mere local
compatibility.
\end{proposition}

The proof is immediate from the fact that the comparison map $\phi$ is canonical: its
construction is symmetric with respect to the roles of $A$ and $B$ up to the
orientation of the overlap, and transitivity follows from the functoriality of the
\v{C}ech complex.

Under Katz-consistency, knowledge integration is governed by cohomological agreement.
Contradictions correspond to nontrivial obstruction classes---to actual failures in
$H^1$---rather than merely to local disagreements. And merging is a higher-categorical
gluing operation that succeeds when the two filtrations of the sections are compatible,
even if the sections themselves are not identical.

\section{The Single Obstruction Engine}

\subsection{The Unified Defect Class}

The preceding sections establish a consistent pattern. In algebraic geometry, in RSVP
field theory, in Spherepop computation, and in semantic infrastructure, the same
structure appears: there are two natural filtrations on the space of admissible
structures, each inducing an obstruction invariant, and the two invariants are always
related by the canonical comparison map $\phi$. This is not a coincidence of notation
but a consequence of a common underlying structure.

Let $\mathcal{F}$ be a prestack or sheaf of admissible structures over a base space
$S$. The failure of global admissibility for a section $s \in \mathcal{F}(S)$ is
measured by an obstruction class
\[
\mathrm{Defect}(s) \in H^1(\mathcal{F}).
\]
This class captures the failure of local data to glue into a global section. It is
independent of the specific method used to probe the system: whether one approaches the
failure differentially or recursively, one finds the same element of $H^1(\mathcal{F})$.

The two filtrations introduced in each domain induce two projections of this class:
\[
\delta_{\mathrm{diff}} \in H^1(\mathrm{gr}_{\mathrm{diff}}(\mathcal{F})),
\qquad
\delta_{\mathrm{rec}} \in H^1(\mathrm{gr}_{\mathrm{rec}}(\mathcal{F})).
\]
The non-abelian Katz formula implies that these projections are related by the
canonical comparison map
\[
\phi : \mathrm{gr}_{\mathrm{rec}}(\mathcal{F}) \to \mathrm{gr}_{\mathrm{diff}}(\mathcal{F}),
\]
such that $\delta_{\mathrm{rec}} = \phi^*(\delta_{\mathrm{diff}})$. Both defects are
images of the same element of $H^1(\mathcal{F})$ under different functorial
realizations, and the comparison map $\phi$ is the natural transformation relating
these two realizations.

\subsection{Symmetry of Failure and Its Implications}

The identification of all admissibility failures with a single class in
$H^1(\mathcal{F})$ yields a symmetry principle: any defect detected through local
infinitesimal analysis must also appear under recursive closure, and vice versa. The
apparent distinction between differential and recursive failure is therefore a
consequence of the chosen representation, not an intrinsic feature of the system.

This has several important implications. First, validation procedures can be unified:
one computes the obstruction class in one regime and transports it via $\phi$, rather
than running independent differential and recursive checks. This reduces the
computational cost of admissibility verification and eliminates the possibility of
inconsistent results from the two checks.

Second, the structure of failure space is stable. Because $H^1(\mathcal{F})$ is a
cohomological group, it is stable under the standard operations of sheaf theory:
restriction, extension, localization, and pushforward. This means that the set of
admissibility failures of a complex system can be computed from the failures of its
components, using the long exact sequence in cohomology. Admissibility is compositional.

Third, the Single Obstruction Engine provides a diagnostic tool. When a system fails,
one can examine the obstruction class $\mathrm{Defect}(s) \in H^1(\mathcal{F})$ to
understand the global topology of the failure. This class is more informative than
either $\delta_{\mathrm{diff}}$ or $\delta_{\mathrm{rec}}$ alone, because it lives
at a higher level of structure and is not tied to any particular representation.

\subsection{Lossy Projection as Structure, Not Deficiency}

The coarse-graining operation by which $p$-curvature is extracted from the full stack
structure might initially appear to be a loss of information: one projects from a
high-fidelity geometric object to a simpler invariant. But the Single Obstruction
Engine reveals that this projection is not a deficiency but a feature.

The full sheared de Rham stack carries all higher coherences. The $p$-curvature is the
obstruction that remains after these coherences are projected away. The remarkable fact
is that this projected invariant is \emph{sufficient}: it captures everything one needs
to know about admissibility failure, because all failures live in the projected
invariant's target space.

This suggests a general principle: in systems governed by Katz-admissibility, a
well-chosen lossy projection does not discard admissibility information. It compresses
the representation without compressing the obstruction. The obstruction is invariant
under the projection; only the higher structure of its witnesses is lost.

In the RSVP context, this means that entropy (the recursive obstruction) is a
complete record of the system's admissibility failures, even though it discards most
of the fine-grained dynamics. In the Spherepop context, the Admissibility Log is a
complete record of the computation's admissibility failures, even though it discards
the intermediate computational states. In the semantic context, the provenance graph
of a knowledge base is a complete record of its mergeability failures, even though it
discards the specific content of each section.

\section{Implications and Future Directions}

\subsection{Physics: Singularities as Comparison-Map Failures}

The identification of differential and recursive obstructions as projections of a
single cohomological class suggests a reformulation of physical law in terms of
admissibility rather than force. In the RSVP framework, the gravitational field and
the entropic structure of a system arise as dual manifestations of the same underlying
defect, related by the linkage equation $\mathcal{E}(X) = \phi^* \mathcal{G}(X)$.

The classical distinction between geometry and thermodynamics is therefore replaced by
a unified theory in which both are expressions of admissibility. This unification is
not achieved by deriving one from the other---entropy is not derived from geometry, nor
geometry from entropy---but by identifying them as projections of one thing onto two
different representational systems.

Singularities receive a new interpretation in this framework. Rather than being points
of divergent curvature or infinite density, they are regions where the comparison map
$\phi$ fails to admit a consistent pullback. At a singularity, the differential
admissibility structure and the recursive admissibility structure diverge so sharply
that no admissible state can be defined there. The singularity is not a point of
physical breakdown; it is a point of representational breakdown, where the two
coordinate systems on admissibility space cease to be compatible.

This interpretation suggests an approach to singularity resolution. Rather than
introducing a cutoff or a regularization, one seeks to extend the comparison map $\phi$
across the singular region by passing to a derived or higher-categorical version of the
admissibility stack. The singularity is resolved when a consistent pullback can be
defined in the extended setting, which may require working with stacks of liftings
rather than individual liftings.

\subsection{Programming Languages: Topological Type Safety}

In computational systems, the replacement of syntactic correctness with topological
agreement leads to a new paradigm for type systems. Programs are trajectories in the
admissibility manifold, and type safety is the condition that these trajectories remain
within the admissible region. Katz-admissibility provides the precise characterization
of what it means to stay within this region: at each step, the differential and
recursive filtrations of the state must be compatible under the comparison map.

This yields a form of type safety that is invariant under reordering, refactoring, or
recomposition, provided the underlying obstruction class remains unchanged. This
invariance is not available in syntactic type systems, where the validity of a program
depends on its syntactic form and can be broken by refactoring even when the semantic
content is preserved.

The Admissibility Log provides a practical mechanism for implementing this invariance.
By maintaining the log as an active filtration of the state-space, the runtime can
compute the recursive defect of any proposed transition without re-examining the entire
history from scratch. The log is a compressed representation of the system's
admissibility history, and Katz-admissibility ensures that this compressed
representation is sufficient for detecting all failures.

Systems of this kind naturally support persistent state, distributed computation, and
incremental verification. The obstruction-theoretic approach to type safety is
particularly well-suited to settings where computation is distributed across multiple
agents, each maintaining its own local state and Admissibility Log, because
Katz-consistency provides a well-defined notion of compatibility between independently
evolving local histories.

\subsection{AI and Knowledge Systems: Learning as Obstruction Minimization}

For artificial intelligence and knowledge representation, the Single Obstruction Engine
provides a mechanism for integrating information without requiring global consensus.
Knowledge bases evolve as collections of partially overlapping sections, each locally
coherent, with admissibility determined by the alignment of obstruction classes under
the Katz-admissibility condition.

This framework supports a reinterpretation of learning. In the standard picture, a
learning system updates its internal representation to minimize a loss function. In
the Katz-admissibility picture, a learning system updates its sections to minimize the
obstruction class $\mathrm{Defect}(s) \in H^1(\mathcal{F})$. The loss function is
replaced by the obstruction class; gradient descent is replaced by a procedure for
moving through the admissibility manifold in the direction of decreasing obstruction.

This reinterpretation has several consequences. First, it provides a topological
characterization of learning: a system has learned successfully when its obstruction
class is trivial, meaning that all local knowledge is globally coherent. Second, it
provides a natural notion of \emph{transfer}: knowledge transfers from one domain to
another when the comparison map $\phi$ can be extended across domain boundaries,
transporting the obstruction class from one context to another. Third, it provides a
characterization of catastrophic forgetting: forgetting occurs when the Admissibility
Log is truncated in a way that introduces a Non-Abelian Tear, breaking the compatibility
between the remaining recursive history and the current differential structure.

Reasoning, in this picture, corresponds to transporting obstruction classes across
different sections of the knowledge sheaf, using the comparison map to verify that the
transport is admissible. An inference is valid when it corresponds to a Katz-admissible
transition in the knowledge state; a logical contradiction corresponds to a
Non-Abelian Tear in the inference graph.

\subsection{Open Problems}

Several directions for further investigation arise naturally from this framework.

The most pressing concerns the explicit computation of the comparison map $\phi$ in
concrete computational and semantic settings. In the algebraic-geometric context, $\phi$
is defined by the relationship between the Hodge and conjugate filtrations on the moduli
stack, and its existence is guaranteed by the non-abelian Katz formula. In
computational and semantic contexts, the filtrations are defined operationally, and the
existence of a canonical comparison map must be verified case by case. A general
existence theorem for Katz-admissibility in categories of computational systems would
be a significant result.

A second open problem concerns the extension of Katz-admissibility to higher-order
obstructions, corresponding to cohomology in degrees greater than one. The present
framework is organized around $H^1(\mathcal{F})$, the first cohomological obstruction.
Higher-degree obstructions, living in $H^n(\mathcal{F})$ for $n \geq 2$, correspond to
subtler forms of incompatibility that cannot be detected by the differential and
recursive projections alone. Extending the Katz formula to these higher degrees would
require working in a derived or $\infty$-categorical setting, where the relevant
comparison maps live in the $\infty$-groupoid of filtrations rather than in an ordinary
category.

A third direction concerns the dynamics of admissibility: how systems evolve within
the admissible region and how they approach its boundary. The present work identifies
the invariant governing admissibility but does not provide a complete description of
the dynamics. A theory of \emph{admissibility flow}---a vector field on the
admissibility manifold that describes how systems move toward or away from
admissibility failure---would complete the framework and provide a basis for predicting
the onset of Non-Abelian Tears before they occur.

Finally, the relationship between Katz-admissibility and the theory of motives deserves
investigation. The $p$-curvature and the Hodge filtration are both aspects of the
crystalline cohomology of a variety, which is related to motivic cohomology. If the
Katz formula can be expressed in motivic terms, it would suggest that admissibility is
a motivic invariant, stable under the operations that define motivic equivalence.

\section{Conclusion}

\subsection{Admissibility as Fundamental Invariant}

The central result of this work is the identification of admissibility as a
cohomological invariant governing systems across geometry, physics, computation, and
semantics. The non-abelian Katz formula, established by Barz \cite{barz}, provides a
rigorous mathematical instance of this principle, demonstrating that differential and
recursive obstructions coincide under a canonical comparison.

This leads to a reformulation of system stability. Rather than requiring consistency in
the sense of logical non-contradiction or convergence to a fixed point, one demands the
alignment of obstruction classes under all admissible transformations. Stability is a
property of the topology of failure rather than the absence of error. A system is
stable not because it never fails but because its failures are coherent: they live in
a single cohomological class whose projections are mutually consistent.

\subsection{The Event Log as Active Constraint}

A significant conceptual contribution of this paper is the reinterpretation of the
event log as an active filtration of the state-space rather than a passive record of
transitions. This reinterpretation has practical consequences. A runtime system that
treats its event log as a second filtration can detect Non-Abelian Tears that classical
type-checking misses. A knowledge system that treats its provenance graph as a
recursive filtration can detect Katz-inadmissible merges that local consistency checks
miss. A physical system in the RSVP framework that treats its entropy as the recursive
projection of the gravitational obstruction can be analyzed with a single unified
formalism rather than separate formalisms for geometry and thermodynamics.

The event log is not what happened. It is a second coordinate system on what is
admissible.

\subsection{Unification Across Domains}

The framework developed here unifies several previously distinct perspectives. In
algebraic geometry, it interprets the relationship between Hodge and conjugate
filtrations as an equivalence of obstruction classes, anchored in Barz's theorem. In
the RSVP framework, it identifies curvature and entropy as dual manifestations of a
single defect related by the linkage equation. In Spherepop, it provides an
operational type rule---the Katz-Admissibility Rule---governing state transitions via
the Admissibility Log. In semantic infrastructure, it yields a criterion for merging
knowledge based on cohomological agreement between differential and recursive
filtrations.

These are not analogies. They are instances of a single mathematical identity
governing systems of different kinds, expressed in different representational
vocabularies. The differences across domains arise only from the choice of
representation, not from the underlying structure. This is the sense in which
Katz-admissibility is a law rather than a framework: it holds independently of
how the systems are described.

\subsection{Final Statement}

We may therefore state the guiding principle in its most concise form:
\[
\text{Admissibility} = \text{Topological Agreement}.
\]
Equivalently: a system is admissible if and only if its differential and recursive
projections of failure coincide under the canonical comparison map $\phi$. A system
fails when and only when these projections diverge, producing a Non-Abelian Tear at
the precise point of divergence. There is no other kind of failure, and no failure that
is visible in one regime but invisible in the other.

\subsection{Outlook}

The identification of a Single Obstruction Engine suggests that many phenomena
traditionally treated as distinct---curvature and entropy, type errors and historical
inconsistency, local incompatibility and provenance conflict---may be unified through
their underlying cohomological structure. Future work may extend this perspective to
higher categories, derived geometries, motivic cohomology, and physical theories beyond
the classical setting.

More broadly, this framework invites a rethinking of how systems are designed, analyzed,
and validated. By shifting the focus from correctness to admissibility, from boolean
checks to cohomological invariants, and from passive records to active filtrations, it
provides a foundation for systems that evolve, interact, and persist while maintaining
coherence across multiple scales of structure. The topology of failure is not an
obstacle to understanding such systems. It is their deepest invariant.

\newpage

\begin{thebibliography}{99}

\bibitem{barz}
M. Barz,
\textit{Non-abelian $p$-curvature and a non-abelian Katz's formula},
arXiv:2604.20054 (2026).

\bibitem{katz}
N. Katz,
\textit{Algebraic solutions of differential equations ($p$-curvature and the Hodge filtration)},
Inventiones Mathematicae, 18 (1972), 1--118.

\bibitem{bhatt}
B. Bhatt,
\textit{$p$-adic Hodge theory notes},
available at \url{https://math.stanford.edu/~bhatt/}.

\bibitem{bhattkanaev}
B. Bhatt, A. Kanaev, A. Mathew, L. Vologodsky, and Z. Zhang,
\textit{Sheared de Rham stacks},
preprint.

\bibitem{drinfeld}
V. Drinfeld,
\textit{On algebraic spaces with an action of $\mathbb{G}_a$},
preprint.

\bibitem{simpson}
C. Simpson,
\textit{Higgs bundles and local systems},
Publications Math\'ematiques de l'IH\'ES, 75 (1992), 5--95.

\bibitem{simpson2}
C. Simpson,
\textit{Moduli of representations of the fundamental group of a smooth projective
variety~I},
Publications Math\'ematiques de l'IH\'ES, 79 (1994), 47--129.

\bibitem{ogus}
A. Ogus,
\textit{$F$-crystals, Griffiths transversality, and the Hodge decomposition},
Ast\'erisque, 221 (1994).

\bibitem{cartier}
P. Cartier,
\textit{Une nouvelle op\'eration sur les formes diff\'erentielles},
C.~R. Acad. Sci. Paris, 244 (1957), 426--428.

\bibitem{illusie}
L. Illusie,
\textit{Complexe cotangent et d\'eformations I},
Lecture Notes in Mathematics, Vol.~239, Springer (1971).

\bibitem{lurie}
J. Lurie,
\textit{Higher Topos Theory},
Annals of Mathematics Studies, Princeton University Press (2009).

\bibitem{maclane}
S. Mac Lane,
\textit{Categories for the Working Mathematician},
Springer (1998).

\bibitem{baez}
J. Baez and J. Dolan,
\textit{Higher-dimensional algebra and topological quantum field theory},
Journal of Mathematical Physics, 36 (1995), 6073--6105.

\bibitem{grothendieck}
A. Grothendieck,
\textit{Crystals and the de Rham cohomology of schemes},
in \textit{Dix Expos\'es sur la Cohomologie des Sch\'emas},
North-Holland (1968), 306--358.

\bibitem{ogusVologodsky}
A. Ogus and V. Vologodsky,
\textit{Nonabelian Hodge theory in characteristic $p$},
Publications Math\'ematiques de l'IH\'ES, 106 (2007), 1--138.

\end{thebibliography}

\end{document}
